\documentclass[12pt]{article}
\usepackage[utf8]{inputenc}
\usepackage[T1]{fontenc}
\usepackage{hyperref}
\title{Southwestern Tarahumara Language Resources for K-12 Curriculum Development}
\author{Compiled by Claude Sonnet 5 (claude-sonnet-5) for Andr\'as Kornai}
\date{}
\begin{document}
\maketitle
\section{Introduction}
Southwestern Tarahumara (ISO 639-3 twr) is an endangered Uto-Aztecan language of Mexico. Searches found almost nothing specific to this variety. Most Tarahumara materials found belong to other varieties and are listed here only as related material.
\section{Primary Educational Resources}
None found for twr.
\section{Linguistic and Reference Materials}
No dictionary or grammar of Southwestern Tarahumara was found. The Joshua Project language page for twr is the only variety-specific item: \url{https://joshuaproject.net/index.php/languages/twr}
\section{Religious and Community Resources}
Joshua Project reports no New Testament for twr. Related varieties (not twr) have New Testaments:
\begin{itemize}
\item Samachique Tarahumara (Central) New Testament: \url{https://dbs.org/bibles/TARWBT}
\item Western (Baja) Tarahumara (tac) New Testament on eBible: \url{https://ebible.org/Bible/details.php?id=tacNT}
\item Scripture Earth page for tac: \url{https://www.scriptureearth.org/00eng.php?iso=tac}
\end{itemize}
Don Burgess's site \url{https://www.tarahumarabooks.com/} offers free downloads, including a Tarahumara/Spanish phrasebook and a Ral\'amuli New Testament. Its variety is not stated.
\section{Digital and Audio Resources}
Audio New Testament for Western (Baja) Tarahumara (tac, not twr): \url{https://megavoice.com/media-cloud/m0b8632-new-testament-tac-western-tarahumara-baja-audio-bible/}
\section{Curriculum Development Context}
Materials for twr itself are essentially absent. Any curriculum would rely on documentation work with older speakers, and on materials from neighbouring varieties only after checking how far they differ from twr.
\end{document}